\documentclass[11pt,a4paper]{article}

\usepackage[utf8]{inputenc}
\usepackage[T1]{fontenc}
\usepackage[brazilian]{babel}
\usepackage{amsmath,amssymb,amsfonts}
\usepackage{booktabs}
\usepackage{geometry}
\usepackage{graphicx}
\usepackage{hyperref}
\usepackage{siunitx}
\usepackage{xcolor}

\geometry{margin=2.5cm}
\hypersetup{colorlinks=true,linkcolor=black,citecolor=black,urlcolor=blue}

\title{Metodologia: Validação da Instabilidade de Saffman--Taylor \\
       em Meio Poroso via Lattice Boltzmann}
\author{Matheus Borges Sampaio}
\date{\today}

\newcommand{\Da}{\mathrm{Da}}
\newcommand{\Ca}{\mathrm{Ca}}
\newcommand{\Bo}{\mathrm{Bo}}
\newcommand{\Cam}{\mathrm{Ca}_m}
\newcommand{\Lref}{L_{\mathrm{ref}}}
\newcommand{\Uref}{U_{\mathrm{ref}}}
\newcommand{\Reff}{\mathrm{Re}}

\begin{document}

\maketitle

\begin{abstract}
Este documento descreve a metodologia adotada para validar a relação de
dispersão linear (\textit{Linear Stability Analysis}, LSA) da instabilidade
de Saffman--Taylor em meio poroso utilizando um simulador baseado em
\textit{Lattice Boltzmann Method} (LBM) acoplado a um modelo de interface
difusa de Cahn--Hilliard. São detalhados (i) o esquema numérico LBM-D2Q9 com
forçamento de Guo, (ii) o modelo de fase, (iii) a extração das amplitudes
interfaciais e da taxa de crescimento exponencial, e (iv) a comparação com a
solução analítica de Darcy. Apresenta-se também a métrica de erro relativo
e três estimadores complementares da taxa de crescimento.
\end{abstract}

% ────────────────────────────────────────────────────────────────────
\section{Visão geral do procedimento}
% ────────────────────────────────────────────────────────────────────

A validação consiste em três etapas:

\begin{enumerate}
    \item \textbf{Simulação numérica:} um deslocamento bifásico em meio
    poroso é inicializado com uma interface perturbada por um modo de
    Fourier $m$ de amplitude $A_0$. O sistema é integrado pelo solver
    LBM--Cahn--Hilliard até $t = T_{\max}$.

    \item \textbf{Extração da taxa numérica $s_{\mathrm{num}}$:} a amplitude
    interfacial $A(t)$ é projetada sobre o modo $m$ via transformada de
    Fourier e ajustada exponencialmente na janela de regime linear,
    fornecendo $A(t) \approx A_0\,e^{s_{\mathrm{num}}\,t}$.

    \item \textbf{Comparação com $s_{\mathrm{ana}}$:} a taxa analítica é
    calculada a partir da relação de dispersão adimensional (Eq.\ \ref{eq:zeta}),
    convertida para unidades de passo de tempo da rede, e o erro relativo
    é definido por $\varepsilon = |s_{\mathrm{num}} - s_{\mathrm{ana}}| /
    |s_{\mathrm{ana}}|$.
\end{enumerate}

% ────────────────────────────────────────────────────────────────────
\section{Modelo LBM--D2Q9 com forçamento de Guo}
% ────────────────────────────────────────────────────────────────────

\subsection{Equação de evolução}

A dinâmica hidrodinâmica é resolvida em um reticulado bidimensional D2Q9
(9 velocidades) sob o operador de colisão BGK com tempo de relaxação único.
Para cada nó $\mathbf{x}$ e velocidade $\mathbf{c}_i$ ($i=0,\ldots,8$):
\begin{equation}
    f_i(\mathbf{x}+\mathbf{c}_i,\,t+1) = f_i(\mathbf{x},t)
    - \frac{1}{\tau}\!\left[f_i - f_i^{\mathrm{eq}}\right]
    + S_i,
    \label{eq:lbm}
\end{equation}
onde $f_i^{\mathrm{eq}}$ é a distribuição de equilíbrio
\begin{equation}
    f_i^{\mathrm{eq}} = w_i \rho \left[ 1 + 3(\mathbf{c}_i\!\cdot\!\mathbf{u})
    + \tfrac{9}{2}(\mathbf{c}_i\!\cdot\!\mathbf{u})^2
    - \tfrac{3}{2}|\mathbf{u}|^2 \right],
\end{equation}
e $S_i$ é o termo de forçamento de Guo, que recupera a equação de
Navier--Stokes com força de corpo até segunda ordem:
\begin{equation}
    S_i = w_i\!\left(1 - \tfrac{1}{2\tau}\right)
    \!\left[ 3(\mathbf{c}_i - \mathbf{u})\!\cdot\!\mathbf{F}_{\mathrm{tot}}
    + 9(\mathbf{c}_i\!\cdot\!\mathbf{u})(\mathbf{c}_i\!\cdot\!\mathbf{F}_{\mathrm{tot}}) \right].
    \label{eq:guo}
\end{equation}

\subsection{Viscosidade efetiva e razão $M$}

O reticulado é monocomponente em $f_i$; as duas fases são distinguidas pelo
campo escalar de fase $\phi(\mathbf{x},t)\in[-1,+1]$, com $\phi=+1$ no
fluido invasor (fase 1) e $\phi=-1$ no fluido deslocado (fase 2). A
viscosidade cinemática local $\nu_{\mathrm{eff}}$ é obtida por mistura
harmônica entre as viscosidades de cada fase:
\begin{equation}
    \frac{1}{\nu_{\mathrm{eff}}} = \frac{S}{\nu_{\mathrm{in}}} + \frac{1-S}{\nu_{\mathrm{out}}},
    \qquad S = \tfrac{1}{2}(\phi+1),
    \label{eq:nu_harm}
\end{equation}
com $\nu_{\mathrm{in}} = (\tau_{\mathrm{in}}-1/2)/3$ e analogamente para
$\nu_{\mathrm{out}}$. A razão de viscosidades $M=\nu_{\mathrm{out}}/\nu_{\mathrm{in}}$
é o parâmetro adimensional dominante da instabilidade. Um modo alternativo
\textit{linear} ($\nu_{\mathrm{eff}} = S\,\nu_{\mathrm{in}} + (1-S)\,\nu_{\mathrm{out}}$)
foi implementado como controle (parâmetro \texttt{VISC\_LINEAR}).

\subsection{Termo de Brinkman/Darcy}

A presença do meio poroso é modelada por uma força volumétrica de
arrasto tipo Darcy:
\begin{equation}
    \mathbf{F}_{\mathrm{Darcy}} = -\sigma_d\,\rho\,\mathbf{u},
    \qquad
    \sigma_d = \frac{1}{K(\mathbf{x})}\!\left( \frac{S}{\nu_{\mathrm{in}}} + \frac{1-S}{\nu_{\mathrm{out}}} \right)^{-1},
\end{equation}
onde $K(\mathbf{x})$ é a permeabilidade absoluta (constante $K_0$ no domínio
homogêneo aqui empregado). A força total que aparece em
(\ref{eq:guo}) é
\begin{equation}
    \mathbf{F}_{\mathrm{tot}} = \mathbf{F}_{\mathrm{cap}}
    + \mathbf{F}_{\mathrm{mag}}
    + \mathbf{F}_{\mathrm{Darcy}}.
\end{equation}
A velocidade física é recuperada com correção implícita do arrasto:
$\mathbf{u} = \mathbf{u}^\star / (1 + \tfrac{1}{2}\sigma_d)$, com
$\mathbf{u}^\star = (\sum_i f_i \mathbf{c}_i + \tfrac{1}{2}\mathbf{F})/\rho$.

% ────────────────────────────────────────────────────────────────────
\section{Modelo de interface: Cahn--Hilliard}
% ────────────────────────────────────────────────────────────────────

\subsection{Energia livre e potencial químico}

A interface difusa é governada pelo modelo Cahn--Hilliard com energia
livre dupla-poço:
\begin{equation}
    \mathcal{F}[\phi] = \int \left[ \beta\,(\phi^2-1)^2
    + \tfrac{\kappa}{2}|\nabla\phi|^2 \right] d\mathbf{x},
\end{equation}
cujo potencial químico é
\begin{equation}
    \mu_\phi = \frac{\delta\mathcal{F}}{\delta\phi}
    = 4\beta\,\phi(\phi^2-1) - \kappa\,\nabla^2\phi.
    \label{eq:mu}
\end{equation}

Para o perfil em equilíbrio $\phi_{\mathrm{eq}}(x) = \tanh(x/W)$, os
parâmetros $\beta$ e $\kappa$ são fixados em função da tensão superficial
$\sigma$ e da espessura nominal $W$:
\begin{equation}
    \beta = \frac{3\sigma}{8W}, \qquad
    \kappa = \frac{3\sigma W}{4}.
    \label{eq:beta_kappa}
\end{equation}
Essa parametrização garante $\sigma = \int_{-\infty}^{\infty}\kappa(\partial_x\phi)^2\,dx$
e foi verificada por testes de Young--Laplace independentes.

\subsection{Evolução temporal e \textit{substepping}}

A equação de Cahn--Hilliard advectiva é
\begin{equation}
    \frac{\partial\phi}{\partial t} + \mathbf{u}\!\cdot\!\nabla\phi
    = M_\phi\,\nabla^2\mu_\phi,
    \label{eq:CH}
\end{equation}
integrada com Euler explícito e operadores de diferenças finitas centradas
de segunda ordem. Devido ao termo $\nabla^4\phi$ contido em $\nabla^2\mu_\phi$,
o passo estável é proporcional a $\Delta x^4$. Para evitar uma redução
agressiva do passo LBM ($\Delta t_{\mathrm{LBM}} = 1$), introduz-se um
\textit{substepping} interno de $N_{\mathrm{sub}} =$ \texttt{CH\_SUBSTEPS}
iterações com $\Delta t_{\mathrm{CH}} = 1/N_{\mathrm{sub}}$:
\begin{equation}
    \phi^{n+1/N_{\mathrm{sub}}} = \phi^n
    + \Delta t_{\mathrm{CH}}
    \left[ M_\phi\,\nabla^2\mu_\phi^n - \mathbf{u}^n\!\cdot\!\nabla\phi^n \right].
\end{equation}
Valores típicos: $M_\phi = 10^{-2}$ a $10^{-3}$,
$N_{\mathrm{sub}} = 60$ a $600$.

\subsection{Força capilar}

A força capilar sobre o reticulado LBM é obtida via formulação
potencial--gradiente:
\begin{equation}
    \mathbf{F}_{\mathrm{cap}} = \mu_\phi\,\nabla\phi,
\end{equation}
que recupera o tensor de tensão de Korteweg
$T_{ij} = -\kappa\,\partial_i\phi\,\partial_j\phi$ no limite contínuo e
gera o salto de pressão correto na interface.

% ────────────────────────────────────────────────────────────────────
\section{Inicialização e condições de contorno}
% ────────────────────────────────────────────────────────────────────

\subsection{Geometria do domínio e perturbação inicial}

O domínio retangular tem dimensões $N_X \times N_Y$ em unidades de rede
(\textit{lattice units}, l.u.). A interface é inicializada como um perfil
$\tanh$ com posição $x_{\mathrm{int}}(y)$ perturbada por um único modo de
Fourier:
\begin{equation}
    x_{\mathrm{int}}(y) = x_c + A_0\,\cos\!\left(\frac{2\pi m y}{N_Y}\right),
    \quad
    \phi_0(x,y) = -\tanh\!\left( \frac{x - x_{\mathrm{int}}(y)}{W} \right),
    \label{eq:phi0}
\end{equation}
onde $m$ é o número de modo (inteiro), $A_0$ a amplitude inicial em l.u., e
$x_c$ a posição média (tipicamente $x_c = 80$). O número de onda físico
adimensional associado é $\alpha = 2\pi m$, tomando $\Lref = N_Y$ como
comprimento de referência.

\subsection{Pré-pressurização de Darcy}

Para evitar o transiente de rampa hidrodinâmica que contaminaria os
primeiros milhares de passos, o campo de densidade $\rho(x,y)$ é
inicializado com o gradiente analítico de Darcy:
\begin{equation}
    \frac{d\rho}{dx} = 3\,\frac{\nu_i(x)}{K_0}\,U_{\mathrm{inlet}},
    \qquad
    i = \begin{cases} \mathrm{in},  & x < x_c \\ \mathrm{out}, & x \geq x_c\end{cases}
\end{equation}
ancorado em $\rho(N_X-1) = 1$, e propagado para trás. Esta inicialização
garante que $\mathbf{u}(t=0) \approx U_{\mathrm{inlet}}\,\hat{\mathbf{x}}$
de maneira consistente com Darcy.

\subsection{Condições de contorno}

\begin{itemize}
    \item \textbf{Entrada} ($x=0$): velocidade prescrita Zou--He com
    $u_x = U_{\mathrm{inlet}}$, $u_y = 0$.
    \item \textbf{Saída} ($x=N_X-1$): pressão prescrita Zou--He com $\rho = 1$.
    \item \textbf{Topo/fundo} ($y=0, N_Y-1$): periódicas em $y$, consistentes
    com a hipótese de modo de Fourier puro.
\end{itemize}

% ────────────────────────────────────────────────────────────────────
\section{Extração da taxa de crescimento numérica}
% ────────────────────────────────────────────────────────────────────

\subsection{Posição da interface}

A cada \textit{snapshot} temporal, lê-se o campo $\phi(x,y,t_n)$ a partir
do arquivo VTK exportado. Para cada linha $y$, localiza-se a posição
$x^\star(y,t_n)$ do contorno $\phi=0$ por interpolação linear entre os
dois nós adjacentes onde ocorre a mudança de sinal:
\begin{equation}
    x^\star(y) = x_\ell + \frac{\phi(x_\ell, y)}{\phi(x_\ell,y) - \phi(x_\ell+1, y)},
\end{equation}
o que fornece precisão sub-pixel ($O(\Delta x^2)$ se $W \gtrsim 3$).

\subsection{Amplitude do modo $m$ via projeção de Fourier}

A amplitude do modo de interesse é extraída por projeção:
\begin{equation}
    A(t_n) = \frac{2}{N_Y}\left| \sum_{y=0}^{N_Y-1}
        x^\star(y, t_n)\,
        \exp\!\left( -\frac{2\pi i\,m\,y}{N_Y} \right) \right|.
    \label{eq:Amp}
\end{equation}
Esta projeção é estritamente mais robusta que a métrica
$A_{\max-\min} = \tfrac{1}{2}(\max_y x^\star - \min_y x^\star)$ porque
\textbf{(i)} filtra modos espúrios de alta frequência, \textbf{(ii)}
suprime contribuições de ruído branco no campo $\phi$ e \textbf{(iii)} é
invariante a translações rígidas da interface.

\subsection{Ajuste exponencial e estimador primário $s_{\mathrm{num}}$}

No regime linear, espera-se $A(t) = A_0\,e^{s\,t}$. Tomando o logaritmo e
realizando regressão linear em mínimos quadrados na janela
$[t_0, t_1]$:
\begin{equation}
    s_{\mathrm{num}} = \mathrm{slope}\!\left[ \log A(t_n) \;\big|\; t_n \in [t_0, t_1] \right].
    \label{eq:snum}
\end{equation}

\subsection{Estimadores complementares: $s_{\mathrm{local}}$}

Para identificar o trecho efetivamente linear da série e detectar viés
do polyfit por inclusão de pontos fora do regime, define-se a taxa
\textbf{instantânea}:
\begin{equation}
    s_{\mathrm{local}}(t_n) = \frac{d}{dt}\log A
    \;\approx\;
    \mathrm{slope}\!\left[ \log A(t_j)\;|\;j\in[n-k,\,n+k] \right],
    \label{eq:sloc}
\end{equation}
calculada por regressão linear deslizante com janela de $2k+1$ pontos
($k$ adaptativo, tipicamente $k = N_{\mathrm{snap}}/16$). A partir dela
extraem-se dois estimadores robustos:

\begin{itemize}
    \item \textbf{Pico}\quad
    $s_{\mathrm{local}}^{\max} = \max_n s_{\mathrm{local}}(t_n)$, melhor
    estimativa pontual da taxa exponencial.

    \item \textbf{Mediana na janela}\quad
    $s_{\mathrm{local}}^{\mathrm{med}} = \mathrm{median}\,\{s_{\mathrm{local}}(t_n)\,:\,t_n\in[t_0,t_1]\}$,
    estimador robusto a curvatura nas bordas da janela.
\end{itemize}

\subsection{Detecção automática da janela de ajuste}

A janela $[t_0,t_1]$ é detectada automaticamente como o platô contíguo de
$s_{\mathrm{local}}(t)$ centrado em seu pico, dentro de uma tolerância
relativa $\delta$:
\begin{equation}
    [t_0, t_1] = \left\{ t_n \;\bigg|\;
    \frac{s_{\mathrm{local}}^{\max} - s_{\mathrm{local}}(t_n)}{s_{\mathrm{local}}^{\max}} \le \delta \right\}_{\text{contíguo a }t_{\mathrm{peak}}}.
\end{equation}
A tolerância é relaxada em escada ($\delta = 4\%, 6\%, 9\%, 13\%, 20\%$)
até a janela conter ao menos $N_{\min} = 6$ pontos.

% ────────────────────────────────────────────────────────────────────
\section{Relação de dispersão analítica (LSA de Darcy)}
% ────────────────────────────────────────────────────────────────────

\subsection{Hipóteses}

A análise linear assume: \textbf{(i)} escoamento de Darcy
em meio poroso homogêneo; \textbf{(ii)} interface abrupta (limite
$W\to 0$); \textbf{(iii)} perturbação de Fourier puro $\delta x(y,t) =
A(t)\cos(ky)$ com $k = 2\pi m/N_Y$; \textbf{(iv)} velocidade base
uniforme $U_{\mathrm{inlet}}$ à montante.

\subsection{Equação 9: forma adimensional}

Adotando $\Lref = N_Y$ e $\Uref = U_{\mathrm{inlet}}$, a taxa adimensional
$\zeta = s\,\Lref/\Uref$ em função do número de onda adimensional
$\alpha = k\,\Lref = 2\pi m$ é
\begin{equation}
    \boxed{\;
    \zeta(\alpha) \;=\;
    \alpha\,\frac{M-1}{M+1}
    \;+\;
    \frac{\Da}{\Ca\,(1+M)}\,\alpha\,
    \Bigl[\, \Bo \;-\; \alpha^2 \;+\; \Cam\,\Lambda_m\,\alpha\,(H_{0n}^2 - H_{0t}^2)\,\Bigr].
    \;}
    \label{eq:zeta}
\end{equation}

Os grupos adimensionais são:
\begin{align}
    M    &= \nu_{\mathrm{out}}/\nu_{\mathrm{in}}
    & \text{(razão de viscosidades)} \\
    \Da  &= K_0/\Lref^2
    & \text{(número de Darcy)} \\
    \Ca  &= \nu_{\mathrm{in}}\,\Uref / \sigma
    & \text{(número capilar)} \\
    \Bo  &= 0
    & \text{(sem gravidade)} \\
    \Cam &= \chi\,H_0^2\,\Lref / \sigma
    & \text{(número de Bond magnético)} \\
    \Lambda_m &= \chi/(2+\chi)
    & \text{(contraste de suscetibilidade)} \\
    H_{0n}^2,\,H_{0t}^2 &= \cos^2\theta,\,\sin^2\theta
    & \text{(componentes do campo)}
\end{align}

A taxa dimensional analítica (em $\mathrm{l.u.}^{-1}$, isto é, por
passo de tempo da rede) é
\begin{equation}
    s_{\mathrm{ana}} = \zeta(\alpha)\,\frac{\Uref}{\Lref}.
    \label{eq:sana}
\end{equation}

\subsection{Limite hidrodinâmico puro}

Nos casos aqui apresentados ($H_0 = 0$, $\chi = 0$, $\Bo = 0$), a Eq.\
(\ref{eq:zeta}) reduz-se a
\begin{equation}
    \zeta(\alpha) = \alpha\,\frac{M-1}{M+1} - \frac{\Da}{\Ca\,(1+M)}\,\alpha^3,
\end{equation}
expressando a competição entre o termo viscoso desestabilizador
$\propto\alpha\,(M-1)/(M+1)$ e o termo capilar estabilizador
$\propto -\alpha^3$. O modo dominante (mais instável) ocorre em
\begin{equation}
    \alpha^\star = \sqrt{ \frac{\Ca\,(M-1)}{3\,\Da} }.
\end{equation}

% ────────────────────────────────────────────────────────────────────
\section{Métrica de erro e protocolo de validação}
% ────────────────────────────────────────────────────────────────────

\subsection{Erro relativo}

Para cada estimador numérico $\hat{s}\in\{s_{\mathrm{num}},\,
s_{\mathrm{local}}^{\max},\,s_{\mathrm{local}}^{\mathrm{med}}\}$, o erro
relativo percentual é definido por
\begin{equation}
    \varepsilon(\hat{s}) = \frac{\bigl|\hat{s} - s_{\mathrm{ana}}\bigr|}{|s_{\mathrm{ana}}|}\times 100\%.
    \label{eq:erro}
\end{equation}

\subsection{Critérios de qualidade da simulação}

A validade da comparação requer que os parâmetros físicos do caso
respeitem os limites do regime linear da LSA. Os critérios monitorados
são:

\begin{enumerate}
    \item \textbf{Resolução da interface:} $W \geq 3$ l.u.\ (critério de
    Nyquist para gradientes em D2Q9). Abaixo disso, correntes espúrias
    contaminam a dinâmica capilar.

    \item \textbf{Razão amplitude/espessura:} $A_0/W \lesssim 0{,}1$ para
    permanecer no regime linear da CH (mantém $|\delta\phi| \ll 1$).
    Valores até $A_0 \sim W/2$ ainda são tolerados na prática.

    \item \textbf{Produto $k\xi$:} a LSA clássica assume interface abrupta;
    o desvio é $\Delta\zeta/\zeta \propto (k\xi)^2$, onde
    $\xi = W\sqrt{2}$ é a espessura efetiva. Valores $k\xi \lesssim 0{,}02$
    mantêm o erro abaixo de 5\%.

    \item \textbf{Conservação de massa:} $|\Delta m|/m_0 < 10^{-3}$ no
    relatório de execução.

    \item \textbf{Concordância entre estimadores:} a diferença
    $|s_{\mathrm{num}} - s_{\mathrm{local}}^{\mathrm{med}}|$ deve ser
    pequena (idealmente menor que o próprio erro vs analítico). Divergência
    significa que a janela de ajuste contém pontos fora do regime
    exponencial puro, e o valor de $s_{\mathrm{num}}$ está viesado.
\end{enumerate}

\subsection{Hierarquia de confiabilidade}

A combinação dos três estimadores fornece informação além de qualquer
um isolado:

\begin{itemize}
    \item \textbf{Concordância dos três} ($|s_{\mathrm{num}}-s_{\mathrm{local}}^{\mathrm{med}}| < 2\,\mathrm{pp}$):
    o regime exponencial foi capturado de forma estável; a taxa medida é
    confiável.

    \item \textbf{Discordância apenas do polyfit:} a janela inclui
    curvatura nas bordas. $s_{\mathrm{local}}^{\mathrm{med}}$ é o
    estimador preferido nesse caso.

    \item \textbf{Discordância dos três:} o caso não está em regime linear
    (amplitude inicial fora dos limites, ou $k\xi$ excessivo). O erro
    reportado não é representativo da física.
\end{itemize}

% ────────────────────────────────────────────────────────────────────
\section{Implementação computacional}
% ────────────────────────────────────────────────────────────────────

O simulador é escrito em Python + Numba (compilação JIT com
\texttt{@njit(parallel=True)}). Os módulos principais são:
\begin{itemize}
    \item \texttt{lbm/lbm.py} --- colisão BGK D2Q9, forçamento de Guo,
    arrasto Darcy, condições Zou--He.
    \item \texttt{cahn\_hilliard/cahn\_hilliard.py} --- evolução temporal
    de $\phi$ com \textit{substepping}.
    \item \texttt{Magnetismo/poisson.py} --- solver SOR para o
    potencial magnético escalar $\widetilde{\psi}$ (desativado nos casos
    deste relatório).
    \item \texttt{initialization/initialization.py} ---
    pré-pressurização de Darcy, perfil $\tanh$ perturbado e equilíbrio
    LBM consistente.
    \item \texttt{post\_process/valida\_lsa.py} --- extração de
    amplitudes, ajuste exponencial e relatório.
\end{itemize}

A integração completa de um caso típico ($N_X\!\times\!N_Y = 920\!\times\!1024$,
$T_{\max} = \num{96000}$ passos) executa em $\sim$\num{2.5} horas em uma
máquina padrão.

% ────────────────────────────────────────────────────────────────────
\section{Resumo do fluxo de validação}
% ────────────────────────────────────────────────────────────────────

\begin{enumerate}
    \item Definir parâmetros do caso em \texttt{casos.json}:
    $N_X$, $N_Y$, $\tau_{\mathrm{in}}$, $\tau_{\mathrm{out}}$,
    $U_{\mathrm{inlet}}$, $K_0$, $\sigma$, $W$, $A_0$, $m$,
    $M_\phi$, $N_{\mathrm{sub}}$.

    \item Executar \texttt{python main.py} $\rightarrow$ geração de
    \texttt{relatorio\_execucao.json} e snapshots VTK.

    \item Executar \texttt{python post\_process/valida\_lsa.py <case\_dir>}
    $\rightarrow$ leitura dos VTKs, extração de $A(t)$ via Eq.\
    (\ref{eq:Amp}), detecção do platô de $s_{\mathrm{local}}$, cálculo
    dos três estimadores e comparação com Eq.\ (\ref{eq:sana}).

    \item Análise: comparar $\varepsilon(s_{\mathrm{num}})$,
    $\varepsilon(s_{\mathrm{local}}^{\max})$ e
    $\varepsilon(s_{\mathrm{local}}^{\mathrm{med}})$ segundo a hierarquia
    de confiabilidade da Seção 6.3.
\end{enumerate}

\end{document}
